\documentclass[11pt,a4paper]{article}
\usepackage[utf8]{inputenc}
\usepackage[vietnamese]{babel}
\usepackage{amsmath,amssymb}
\usepackage{geometry}
\geometry{a4paper, total={180mm,260mm}, left=15mm, right=15mm, top=18mm, bottom=18mm}
\usepackage{tikz}
\usetikzlibrary{shapes.geometric, arrows.meta, positioning, calc}
\usepackage{tcolorbox}
\tcbuselibrary{skins,breakable}
\usepackage{xcolor}
\usepackage{multicol}
\usepackage{array}
\usepackage{fontawesome}
\usepackage{fancyhdr}

\pagestyle{fancy}
\fancyhf{}
\renewcommand{\headrulewidth}{0.5pt}
\renewcommand{\footrulewidth}{0.5pt}
\fancyhead[L]{\small\bfseries\color{teal}\faBook\ SỔ TAY KIẾN THỨC TOÁN 3}
\fancyhead[R]{\small\bfseries\color{teal}Lớp Toán Thầy Lĩnh -- loptoanthaylinh.vn}
\fancyfoot[C]{\small\color{gray}-- Sổ tay kiến thức: Bài 51. Diện tích của một hình. Xăng-ti-mét vuông ($\text{cm}^2$) --}

\definecolor{primary}{HTML}{005F73}
\definecolor{secondary}{HTML}{0A9396}
\definecolor{accent}{HTML}{EE9B00}
\definecolor{danger}{HTML}{AE2012}
\definecolor{boxbg}{HTML}{F4F6F8}

\begin{document}

% --- BANNER TIÊU ĐỀ ---
\begin{tcolorbox}[
  enhanced,
  colback=white,
  colframe=primary,
  arc=3mm,
  boxrule=1.2pt,
  left=4mm, right=4mm, top=2.5mm, bottom=2.5mm,
  drop shadow=black!10!white
]
  \centering
  {\Large\bfseries\color{primary}\faBookmark\ BÀI 51: DIỆN TÍCH CỦA MỘT HÌNH. XĂNG-TI-MÉT VUÔNG ($\text{cm}^2$)}\\[1.5mm]
  {\small\itshape\color{black!75}Khái niệm diện tích bề mặt, đơn vị đo xăng-ti-mét vuông, phương pháp đếm và ghép ô vuông trên 1 trang A4}
\end{tcolorbox}

\vspace{-2mm}


\noindent
\begin{minipage}[t]{0.49\textwidth}\vspace{0pt}
\begin{tcolorbox}[
  enhanced,
  colback=blue!4!white,
  colframe=blue!75!black,
  arc=2mm, boxrule=0.8pt,
  left=2mm, right=2mm, top=1.8mm, bottom=1.8mm,
  title={\footnotesize\bfseries\color{white}\faKey\ 1. Khái niệm Diện tích của một hình}
]
{\footnotesize
\begin{itemize}\setlength{\itemsep}{1.5pt}\setlength{\leftskip}{-2.5mm}
  \item Diện tích của một hình là \textbf{toàn bộ bề mặt} phần mặt phẳng nằm trọn bên trong đường bao của hình đó.
  \item \textbf{So sánh diện tích 2 hình:}
  \begin{itemize}\setlength{\itemsep}{0.5pt}\setlength{\leftskip}{-2.5mm}
    \item Hình $A$ nằm hoàn toàn trong hình $B$ $\implies$ Diện tích hình $A$ nhỏ hơn diện tích hình $B$.
    \item Hai hình có cùng số ô vuông đơn vị $\implies$ Diện tích bằng nhau (dù hình dạng khác nhau).
    \item Tách hình: Diện tích hình lớn bằng \textbf{tổng diện tích} các hình thành phần được ghép lại.
  \end{itemize}
\end{itemize}
}
\end{tcolorbox}
\end{minipage}
\hfill
\begin{minipage}[t]{0.49\textwidth}\vspace{0pt}
\begin{tcolorbox}[
  enhanced,
  colback=teal!4!white,
  colframe=teal!75!black,
  arc=2mm, boxrule=0.8pt,
  left=2mm, right=2mm, top=1.8mm, bottom=1.8mm,
  title={\footnotesize\bfseries\color{white}\faSliders\ 2. Đơn vị Xăng-ti-mét vuông ($\\text{cm}^2$)}
]
{\footnotesize
\begin{itemize}\setlength{\itemsep}{1.5pt}\setlength{\leftskip}{-2.5mm}
  \item Xăng-ti-mét vuông là diện tích của hình vuông có cạnh dài đúng $\mathbf{1\text{ cm}}$.
  \item Xăng-ti-mét vuông viết tắt là $\mathbf{\text{cm}^2}$ (chữ $\text{cm}$ kèm số $2$ nhỏ ở phía trên bên phải).
  \item \textbf{Thực hiện phép tính có $\text{cm}^2$:}
  \begin{itemize}\setlength{\itemsep}{0.5pt}\setlength{\leftskip}{-2.5mm}
    \item $15\text{ cm}^2 + 25\text{ cm}^2 = \mathbf{40\text{ cm}^2}$.
    \item $12\text{ cm}^2 \times 3 = \mathbf{36\text{ cm}^2}$.
    \item $48\text{ cm}^2 : 6 = \mathbf{8\text{ cm}^2}$.
  \end{itemize}
\end{itemize}
}
\end{tcolorbox}
\end{minipage}

\vspace{-2.5mm}
% --- KHỐI 3: SƠ ĐỒ TƯ DUY TIKZ ---
\noindent
\begin{tcolorbox}[
  enhanced,
  colback=orange!4!white,
  colframe=orange!80!black,
  arc=2mm, boxrule=0.8pt,
  left=2mm, right=2mm, top=1.5mm, bottom=1.5mm,
  title={\footnotesize\bfseries\color{white}\faSitemap\ SƠ ĐỒ TƯ DUY: HÌNH THÀNH KHÁI NIỆM DIỆN TÍCH VÀ ĐƠN VỊ ĐO}
]
\centering
\begin{tikzpicture}[
  box/.style={draw, rounded corners=2pt, align=center, font=\scriptsize\bfseries, inner sep=2.5pt, minimum height=0.75cm},
  arr/.style={-{Stealth[length=2mm]}, thick}
]
  \node[box, fill=purple!15, draw=purple!80!black, text width=2.6cm] (s1) at (0, 0) {BỀ MẶT CỦA HÌNH\\Toàn bộ phần phẳng\\nằm bên trong hình};
  \node[box, fill=teal!15, draw=teal!80!black, text width=2.8cm] (s2) at (3.8, 0) {ĐƠN VỊ ĐO $\text{cm}^2$\\Hình vuông cạnh $1\text{ cm}$\\Kí hiệu: $\mathbf{\text{cm}^2}$};
  \node[box, fill=blue!12, draw=blue!80!black, text width=2.8cm] (s3) at (7.6, 0) {ĐẾM Ô VUÔNG\\Số ô vuông $1\text{ cm}^2$\\$\implies$ Diện tích hình};
  \node[box, fill=orange!15, draw=orange!80!black, text width=2.4cm] (s4) at (11.0, 0) {TÍNH TOÁN\\Cộng, trừ, nhân, chia\\kèm $\text{cm}^2$};

  \draw[arr, purple!80!black] (s1) -- (s2);
  \draw[arr, teal!80!black] (s2) -- (s3);
  \draw[arr, blue!80!black] (s3) -- (s4);
\end{tikzpicture}
\end{tcolorbox}

\vspace{-2.5mm}
% --- KHỐI 4: VÍ DỤ & MẸO TRÁNH BẪY ---
\noindent
\begin{minipage}[t]{0.49\textwidth}\vspace{0pt}
\begin{tcolorbox}[
  enhanced,
  colback=green!4!white,
  colframe=green!60!black,
  arc=2mm, boxrule=0.8pt,
  left=2mm, right=2mm, top=1.8mm, bottom=1.8mm,
  title={\footnotesize\bfseries\color{white}\faCheckCircle\ Ví dụ thực hành chuẩn mực}
]
{\footnotesize
\textbf{1. Đọc và viết đơn vị:}\\
-- $250\text{ cm}^2$ đọc là: \textit{Hai trăm năm mươi xăng-ti-mét vuông}.\\
-- \textit{Một nghìn không trăm lẻ năm xăng-ti-mét vuông} viết là: $\mathbf{1005\text{ cm}^2}$.\\
\textbf{2. Bài toán tách và ghép hình:}\\
Hình $H$ gồm $6$ ô vuông, mỗi ô vuông có diện tích $1\text{ cm}^2$. Diện tích hình $H$ là:\\
$1 \times 6 = \mathbf{6\text{ cm}^2}$.
}
\end{tcolorbox}
\end{minipage}
\hfill
\begin{minipage}[t]{0.49\textwidth}\vspace{0pt}
\begin{tcolorbox}[
  enhanced,
  colback=red!4!white,
  colframe=red!70!black,
  arc=2mm, boxrule=0.8pt,
  left=2mm, right=2mm, top=1.8mm, bottom=1.8mm,
  title={\footnotesize\bfseries\color{white}\faExclamationTriangle\ Mẹo tránh bẫy trắc nghiệm}
]
{\footnotesize
\begin{itemize}\setlength{\itemsep}{1.5pt}\setlength{\leftskip}{-2.5mm}
  \item \textbf{Bẫy nhầm giữa cm (độ dài) và $\text{cm}^2$ (diện tích):}\\
  Chu vi là độ dài đường bao (đo bằng $\text{cm}$). Diện tích là bề mặt (đo bằng $\text{cm}^2$). Nhầm đơn vị là sai bản chất hình học!
  \item \textbf{Bẫy diện tích bằng nhau thì chu vi phải bằng nhau:}\\
  Hai hình có cùng diện tích có thể có chu vi hoàn toàn khác nhau (ví dụ chữ nhật $1 \times 4$ và $2 \times 2$).
\end{itemize}
}
\end{tcolorbox}
\end{minipage}


\end{document}
